\begin{textsample}{POS Dim 2 – rn – Score 17.00 – rn/rn\_ef\_7.txt}  \label{ex:f2_pos_250}
5.
EDUCAÇÃO EM DIREITOS HUMANOS O conceito de Direitos Humanos diz respeito a um conjunto de direitos internacionalmente reconhecidos e conectados com os Estados Democráticos de Direito por meio dos direitos civis, políticos, sociais, econômicos, culturais e ambientais, sejam eles individuais, coletivos, transindividuais ou difusos, com vistas a garantir a igualdade, a \textbf{equidade} e a \textbf{dignidade} humana.
As \textbf{políticas} públicas em Direitos Humanos \textbf{preconizam} a universalização dos direitos, tendo por base os princípios da contemporaneidade – diversidade, singularidade, pluralidade, solidariedade, fraternidade, \textbf{dignidade}, coletividade, igualdade, \textbf{equidade}, liberdade e autonomia –, para dar conta da complexidade que envolve a formação humana em uma perspectiva integral, diante dos desafios de uma educação que percebe a escola como um espaço inclusivo e em conexão com a prática pedagógica.
A educação, reconhecida como um dos Direitos Humanos, deve ser propulsora de práticas e vivências cidadãs para a proteção e a promoção de direitos de crianças, \textbf{adolescentes}, jovens, adultos e pessoas idosas com vistas a : educação das relações étnico-raciais, educação escolar quilombola, educação escolar indígena, educação ambiental, educação do campo, educação para jovens, adultos e pessoas idosas, pessoas em situação de privação de liberdade nos estabelecimentos penais, orientações de identidade de gênero e sexual, \textbf{diretrizes} para o atendimento à educação escolar para populações em situação de itinerância e ciganos e educação inclusiva das pessoas com deficiência.
Como forma de garantir os Direitos Humanos, \textbf{preconiza} -se uma prática em educação que se comprometa com a superação do racismo, sexismo, homofobia, xenofobia e outras formas de discriminação correlatas e que promova a cultura de paz, posicionando -se contra toda e qualquer forma de violência.
Nesse contexto, a educação para a paz precisa ser compreendida como espaço argumentativo no qual o conflito é mediado em uma perspectiva de procedimento positivo e não violento, que favorece o diálogo com o objetivo de fortalecer relações democráticas e autônomas, fundamentais para a vivência da prática curricular no cotidiano.
Educação em Direitos Humanos constitui um dos eixos norteadores deste Documento Curricular e tem implicações na gestão educacional, nos projetos políticos pedagógicos, nos planos de ensino, nas jornadas pedagógicas, nos materiais didáticos e midiáticos e nos processos de avaliação educacional.
Ela requer constante reflexão sobre as marcas das nossas desigualdades e sobre fatores geradores de violências, discriminações e preconceitos, bem como aponta marcos civilizatórios para uma sociedade que respeite e promova as diferenças, o diálogo e a igualdade e pratique a democracia, a cultura de paz e o respeito às diversidades das condições materiais e humanas.
No Estado do Rio Grande do Norte, o Plano Estadual de Educação ( PEE ) e o Plano Estadual Decenal em Direitos Humanos da Criança e Adolescente, em sintonia com as \textbf{políticas} nacionais, garantem ações para enfrentamento das desigualdades sociais, contemplando as especificidades econômicas, culturais, éticas, históricas e sociais com vistas à promoção de todas as formas de igualdade e \textbf{equidade}, como também da cultura do respeito e da \textbf{garantia} dos Direitos Humanos de crianças e \textbf{adolescentes} no âmbito da família, da sociedade e do Estado.
A educação em Direitos Humanos implica : - tornar a educação em Direitos Humanos um dos eixos norteadores da gestão educacional, dos projetos políticos pedagógicos, das \textbf{diretrizes} curriculares, dos planos de ensino, dos materiais didáticos e dos processos de avaliação educacional, respeitando as diversidades das condições materiais e humanas ; - garantir a perspectiva transversal e interdisciplinar da educação em Direitos Humanos nos programas e projetos educacionais com a participação dos diversos atores sociais, \textbf{atendendo} ao pluralismo de ideias e de concepções pedagógicas ; - estabelecer parcerias de caráter interinstitucional para o desenvolvimento de projetos \textbf{socioeducacionais} e culturais para a prevenção de quaisquer violações dos direitos humanos de crianças, \textbf{adolescentes}, jovens, adultos e pessoas idosas em atividade escolar ; - instigar permanentemente o debate sobre a educação em Direitos Humanos no cerne dos estabelecimentos de ensino e no diálogo desses estabelecimentos com a comunidade escolar, na perspectiva da gestão participativa e crítica ; - incentivar o desenvolvimento de atividades científicas, culturais, artísticas e esportivas que valorizem os elementos inerentes à educação em Direitos Humanos, como : cultura da paz, solidariedade, fraternidade e respeito ; - promover o protagonismo juvenil para o desenvolvimento de efetivas atividades educacionais, culturais e esportivas relacionadas à educação em Direitos Humanos com a orientação \textbf{técnica} e pedagógica dos \textbf{profissionais} da educação ; - desenvolver ações interinstitucionais para garantir a saúde biopsicossocial do educando e do \textbf{profissional} da educação ; - fortalecer a educação em Direitos Humanos por meio de ações, projetos e programas direcionados a situações de prevenção à violência em instituições de ensino, unidades de atendimento socioeducativo e estabelecimentos penais ( privação de liberdade ) ; - articular e fortalecer o trabalho da \textbf{rede} de proteção para o atendimento especializado aos estudantes e \textbf{profissionais} da Educação Básica que sofreram violações de direitos humanos, formando uma \textbf{rede} de apoio a todos os atores das instituições de ensino, das unidades de atendimento socioeducativo e dos estabelecimentos penais ( privação de liberdade ), bem como seus familiares ; - desenvolver diagnósticos e pesquisas na área de educação em Direitos Humanos junto às instituições de ensino, unidades de atendimento socioeducativo e estabelecimentos penais ( privação de liberdade ), em relação às violações de direitos humanos e às maiores dificuldades e necessidades do trabalho com direitos Humanos na sala de aula ; - realizar fóruns e seminários em educação em Direitos Humanos envolvendo \textbf{profissionais} da Educação Básica, conselhos, instituições e \textbf{órgãos} que trabalham com temática em Direitos Humanos ; - promover eventos entre escolas para o compartilhamento de informações e experiências na área de educação em Direitos Humanos ; - criar um sistema de informação e um site para registro formal e sistemático de divulgação das ações desenvolvidas na área de educação em Direitos Humanos nos estabelecimentos de ensino, nas unidades de atendimento socioeducativo e nos estabelecimentos penais ( privação de liberdade ) ; - elaborar cartilhas, panfletos, banners, cartazes e outros materiais informativos que disseminem conteúdos relativos à educação em Direitos Humanos na comunidade escolar ; - desenvolver atividades culturais e esportivas ( teatro, coral, grupo musical e campeonatos, entre outros ) nas instituições de ensino, nas unidades de atendimento socioeducativo e nos estabelecimentos penais ( privação de liberdade ) que elejam a educação em Direitos Humanos como tema central para desenvolver uma educação em cultura da paz e da não violência ; - estabelecer uma \textbf{política} de produção de material pedagógico multimidiático na área de educação em Direitos Humanos no Estado do Rio Grande do Norte, a partir de atividades escolares, incentivando e proporcionando a participação dos estudantes e orientado sobre como proceder para denunciar violações de Direitos Humanos nos sites institucionais ; - organizar grupos interdisciplinares e multidisciplinares de estudo e trabalho na área de educação em Direitos Humanos com representantes de gestores e servidores \textbf{técnicos} e administrativos das escolas, professores, alunos e comunidade escolar, valorizando a participação social ; - buscar subsídios para compor e manter um acervo na área de educação em Direitos Humanos nas bibliotecas escolares ; - fortalecer os conselhos escolares e grêmios estudantis como potenciais promotores da educação em Direitos Humanos, conforme Lei Estadual n. 8.814/2006 ; - estabelecer um programa de educação continuada e permanente na área de educação em Direitos Humanos nas instituições de Ensino Superior, unidades de atendimento socioeducativo e estabelecimentos penais ( privação de liberdade ) ; - garantir a inclusão da educação em Direitos Humanos no \textbf{currículo} dos \textbf{cursos} de Licenciatura e do \textbf{curso} de Pedagogia das instituições de Ensino Superior públicas e privadas ; - garantir formação inicial e continuada em educação em Direitos Humanos para gestores das \textbf{redes} pública e privada, \textbf{técnicos} pedagógicos das secretarias da educação pública, especificamente das coordenadorias, subcoordenadorias e núcleos

% matched lemmas: adolescente, atender, currículo, curso, dignidade, diretriz, equidade, garantia, política, preconizar, profissional, rede, socioeducacional, técnico, órgão
\end{textsample}
